\section{Week 8}
\sol{8.1}{Let Ada accept tasks $1,2$ and Ben accept only task $1$. The matching consisting of Ada--$1$ is maximal: Ben--$1$ shares the used task and Ada--$2$ shares the used student, so neither can simply be added. It is not maximum because the matching Ada--$2$, Ben--$1$ has size two. Maximal means no edge can be added without revising the current choices; maximum compares against every possible matching.}
\sol{8.2}{The unmatched student is $c$. Use the alternating path
\[
 c\ --\ 2\ --\ a\ --\ 1\ --\ b\ --\ 3.
\]
The edges $c$--$2$, $a$--$1$, and $b$--$3$ are outside the matching; $a$--$2$ and $b$--$1$ are inside it. Task $3$ is unmatched. Flipping the path gives $c$--$2$, $a$--$1$, $b$--$3$, a matching of size three. Each internal vertex exchanges one matched edge for another, while the endpoints become matched.}
\sol{8.3}{Let $S$ be all three students. Their neighbor set is $N(S)=\{1,2\}$, so $|N(S)|=2<3=|S|$. Distinct assignments of all three would require three distinct members of a two-element set, which is impossible. Counting each student's two options separately counts the same tasks repeatedly; feasibility depends on the number of distinct options available to the group.}
\sol{8.4}{The flip proof treats each internal vertex as incident to exactly two path edges, one matched and one unmatched. If a vertex appeared twice, it could meet two unmatched path edges, and flipping would then give it two partners, so the result would not be a matching.

To avoid repeated vertices, think of the search as moving along arrows: an unmatched edge may be used only from a student to a task, and a matched edge only from a task to its student. Call the resulting directed graph the \emph{alternating-search graph}. Every directed route in it from the unmatched root alternates correctly, because the arrows only allow the permitted moves. Take a shortest such route from the root to an unused task. If it visited some vertex twice, deleting the segment between the two visits would leave a shorter directed route with the same endpoints, just as for walks in Chapter~7, and it would still alternate because it still follows the arrows. That contradicts shortestness. So a shortest route has no repeated vertex and is an augmenting path. The predecessor records kept by the search produce such a path directly, because each vertex is recorded only once.}
\sol{8.5}{Increasing a quantity by one at every step proves termination only if the quantity cannot grow forever: the natural numbers $0,1,2,\ldots$ increase by one at each step and never stop. Finiteness of $L$ supplies the needed ceiling. A matching uses each left vertex at most once, so its size is at most $|L|$. Starting with size $k$, and increasing it by one at each augmentation, allows at most $|L|-k$ augmentations. From the empty matching the bound is $|L|$. Without a finite upper bound, a quantity increasing by one can increase indefinitely; strict progress alone is not a termination proof.}
\sol{8.6}{The search starts at one unmatched left vertex $u$; other matched and unmatched vertices are explored by the rules in Chapter~8. Every reached right vertex was reached from a reached left vertex along an edge, so $B\subseteq N(A)$. For the converse, let $a\in A$ and let $b$ be any neighbor of $a$. If $ab$ is unmatched, the search explores that left-to-right edge and reaches $b$. If $ab$ is matched, then $a\ne u$. The only way the search reaches this other left vertex is through its unique matched right partner, which must be $b$. Thus $b$ was already reached. In either case $b\in B$, proving $N(A)\subseteq B$. Exhaustion of the search is needed to say all eligible unmatched edges have been explored.}
